%   % !TEX root = ../../VIII,3_Rahmen-TeX_8-1.tex
%
%
%   Band VIII, 3 N.~??A16 // AEF
%   Signatur/Tex-Datei: LH_35_09_15_020r
%   RK-Nr. 60650 (= ex 41154 B; ursprüngliche Zuschreibung korrigiert: RK 41154 und 41155 gehören nicht zusammen)
%   Überschrift: [De corporum compositorum duritie]
%   Datierung: [Ende Januar 1683 -- 1684 = RK 60651]
%   WZ: LEd-WZ 803006, Gegenmarke G D (= RK-WZ: 273)
%.  SZ: (keins)
%.  Bilddateien (PDF): LH_35_09_15_020r_d (insgesamt eins)
%
%
\selectlanguage{ngerman}%
\frenchspacing%
%
\begin{ledgroupsized}[r]{120mm}
\footnotesize
\pstart
\noindent\textbf{Überlieferung:}
\pend
\end{ledgroupsized}
\begin{ledgroupsized}[r]{114mm}
\footnotesize
\pstart \parindent -6mm
\makebox[6mm][l]{\textit{L}}%
Notiz: LH~XXXV~9,~15 Bl.~20.
Ein Blatt~4\textsuperscript{o};
Fragment eines Wasserzeichens am Blatt\-rand:
Papier aus dem Harz.
Eine Seite auf Bl.~20~r\textsuperscript{o};
Bl.~20~v\textsuperscript{o} überliefert N.~18.
\pend
\end{ledgroupsized}
%
%
\vspace*{5mm}
\begin{ledgroup}
\footnotesize
\pstart
\noindent%
\textbf{Datierungsgründe:}
Die vorliegende Notiz N.~19 ist auf demselben Textträger wie die Aufzeichnung N.~18 überliefert.
Diese letztere lässt sich auf den Zeitraum von Ende Januar 1683 bis Ende des Jahres 1684 datieren (siehe die entsprechende Begründung, S.~\pageref{LH_35_09_15_020v_Datierung}).
Aufgrund der gemeinsamen Überlieferung liegt es nahe, für N.~19 dieselbe Datierung vorzuschlagen.
\pend
\end{ledgroup}
%
%
\selectlanguage{latin}%
\frenchspacing%
%
%
\vspace*{8mm}
\pstart%
\normalsize%
\noindent%
%
\lbrack20~r\textsuperscript{o}\rbrack\ %%%% Blatt 20r
%
Ex cupro\protect\index{Sachverzeichnis}{cuprum} et stanno\protect\index{Sachverzeichnis}{stannum} confusis,
fit corpus durissimum,\protect\index{Sachverzeichnis}{corpus durum}
magis quam composita,\protect\index{Sachverzeichnis}{corpus compositum}
et in talibus casibus\protect\index{Sachverzeichnis}{casus} non est
\edtext{verum, non dare}{%
\lemma{verum,}\Bfootnote{%
\hspace{-0,5mm}\textbar~non \textit{erg.}~%
\textbar\ dare%
~\textit{L}}}
quid posse,
quod non habet.
% [[\textit{abbrechend??}]]
\pend%
%
\pstart%
\edtext{Ex glutine\protect\index{Sachverzeichnis}{gluten} et}{%
\lemma{Ex}\Bfootnote{%
\textit{(1)}~colla\protect\index{Sachverzeichnis}{colla} et
\textit{(2)}~glutine et%
~\textit{L}}}
charta\protect\index{Sachverzeichnis}{charta} fit etiam corpus
durius\protect\index{Sachverzeichnis}{corpus durus} longe et charta et glutine.
Gluten habet duritiem,\protect\index{Sachverzeichnis}{durities}
sed est fragile,\protect\index{Sachverzeichnis}{fragilitas}
charta habet mollitiem,\protect\index{Sachverzeichnis}{mollities}
et est tensilis:\protect\index{Sachverzeichnis}{tensilitas}
si misceantur inter se non potest glutinis unius pars ab alia abjungi,
nisi simul chartae pars una ab alia sejungatur,
ergo minor
\edtext{est distensio\protect\index{Sachverzeichnis}{distensio} tam}{%
\lemma{est}\Bfootnote{%
\textit{(1)}~tam
\textit{(2)}~distensio tam%
~\textit{L}}}
glutinis quam chartae.
\pend%
\vspace{0.5em}%
\pstart%
\noindent%
\lbrack\textit{Nachfolgend kleingedruckter Text in L gestrichen:}\rbrack%
\pend%
\footnotesize%
\vspace{0.5em}%
\pstart%
\noindent%
At notandum cum
\edtext{glutinis frustum\protect\index{Sachverzeichnis}{frustum} \textit{A} frangimus}{%
\lemma{glutinis}\Bfootnote{%
\textit{(1)}~partem fran
\textit{(2)}~frustum \textit{A} frangimus%
~\textit{L}}}
nos operari quasi per vectem\protect\index{Sachverzeichnis}{vectis}
cuius fulcrum\protect\index{Sachverzeichnis}{fulcrum} \textit{B}.
Et ita minore vi\protect\index{Sachverzeichnis}{vis frangendi} opus habere quam recta diducendo.
\pend%
%
  \vspace{1.5em}%
  \centerline{\includegraphics[width=0.12\textwidth]{gesamttex/edit_VIII,3/images/LH_35_09_15_020r_d.pdf}}%
  \vspace*{0.5em}
  \centerline{{\normalsize{\lbrack\textit{Fig.~1, gestr.}\rbrack}}}%
  \vspace{1.5em}%
%
\normalsize%
\pstart%
\noindent%
Charta autem per se ferre potest multo majorem distensionem\protect\index{Sachverzeichnis}{distensio}
quam gluten per se,
contra gluten minus cedit quam charta.
Charta non facile frangitur, gluten non facile flectitur.
Ergo compositum non facile flectitur, nec facile frangitur.
Sed quaeritur
an compositum\protect\index{Sachverzeichnis}{corpus compositum}
difficilius frangatur quam charta,\protect\index{Sachverzeichnis}{charta}
et difficilius flectatur quam gluten.\protect\index{Sachverzeichnis}{gluten}%
%
%% ! ! ! !    ACHTUNG GETRIXT: Folgende Cfootnote hängt eigentlich an Fig. 1    ! ! ! ! 
%\edtext{}{\lemma{\hspace{1,6mm}\lbrack\textit{Fig.~1}\rbrack}\killnumber\Cfootnote{%
%Siehe den textkritischen Apparat zu Z. ??12.}}
\pend%
%%
%% \newpage
% \vspace*{2.0em}%
%%
%  \centerline{\includegraphics[width=0.16\textwidth]{gesamttex/edit_VIII,3/images/LH_35_09_15_020r_d.pdf}}%
%  \vspace*{0.5em}
%  \centerline{\lbrack\textit{Fig.~1, gestr.}\rbrack}%
%%
%% \newpage
% \vspace*{1.0em}%
%
%
% ENDE DES STÜCKES auf Blatt 20r